% generated by roam.py handbook export -- не редактировать, править заметки
% 20260916143716-выпуклость_проекции.org (body ba73a7f085ce)

\begin{theorem}[{Выпуклость проекции}]\label[theorem]{note:80AFB48E-12EF-4D6C-9303-E2E52FC83D26}
  Пусть:

  \begin{itemize}
    \item $V$~--- \hyperref[note:5D688904-8445-47B5-9D39-8BCC39F9C354]{гильбертово пространство},
    \item $U \subset V$~--- замкнутое подпространство,
    \item $C \subset V$~--- \hyperref[note:316B5BFA-8B9F-4267-9086-D6593AAB5F9B]{выпуклое множество}
    \item $P_U: V \rightarrow U$~--- проекция (по свойствам гильбертовости это линейный оператор)
  \end{itemize}

  Тогда: $P_U C$~--- тоже выпуклое
\end{theorem}

\begin{proof}
  $x, y \in P_UC$, следовательно у них есть прообразы: $x = P_Ua, y = P_Ub$, где $a,b \in C$.
  Мы хотим доказать, что $\forall \lambda \in [0;1]\ \forall x, y \in P_UC\ \lambda x + (1-\lambda)y \in P_UC$.
  Заменим $x,y$ на проекции прообразов: $\lambda x + (1-\lambda) y = \lambda P_U a + (1- \lambda) P_U b$.
  По линейности: $\lambda P_Ua + (1-\lambda)P_U b = P_U(\lambda a + (1-\lambda) b)$.
  Но мы знаем, что $\lambda a + (1-\lambda)b \in C$ (выпуклость), тогда
  $P_U(\lambda a + (1-\lambda) b) = \lambda P_U a + (1-\lambda) P_U b \in P_U C$
\end{proof}
